\documentclass[10pt,a4paper]{article}

% =========================================================
% PACKAGES
% =========================================================
\usepackage[utf8]{inputenc}
\usepackage[T1]{fontenc}
\usepackage[french]{babel}
\usepackage{lmodern}
\usepackage[
    a4paper,
    top=0.55cm,
    bottom=0.5cm,
    left=0.85cm,
    right=0.85cm
]{geometry}
\usepackage{enumitem}
\usepackage{hyperref}
\usepackage{titlesec}
\usepackage{microtype}

% =========================================================
% HYPERLINKS
% =========================================================
\hypersetup{
    colorlinks=true,
    urlcolor=black,
    linkcolor=black,
    pdfborder={0 0 0}
}

% =========================================================
% GENERAL SETTINGS
% =========================================================
\setlength{\parindent}{0pt}
\setlength{\parskip}{0pt}
\setlength{\baselineskip}{9.5pt}

\setlist[itemize]{
    leftmargin=1.1em,
    itemsep=0pt,
    topsep=0pt,
    parsep=0pt,
    partopsep=0pt
}

% =========================================================
% SECTION STYLE
% =========================================================
\titleformat{\section}
    {\normalsize\bfseries}
    {}
    {0pt}
    {}
    [\titlerule]

\titlespacing{\section}
    {0pt}
    {2.5pt}
    {1pt}

% =========================================================
% DOCUMENT
% =========================================================
\begin{document}

% =========================================================
% HEADER
% =========================================================

\begin{center}

{\Large \textbf{TAKWA LAFFET}}

\vspace{0pt}

{\small
\textbf{Ingénieure en Cybersécurité | Full-Stack Developer | DevSecOps}
}

\vspace{1pt}

Tunisie \;|\;
\textbf{+216 93 250 946} \;|\;
\href{mailto:takwa.laffet@esprit.tn}{takwa.laffet@esprit.tn}

\vspace{1pt}

\href{https://www.linkedin.com/in/takwa-laffet-883239211/}
{LinkedIn: linkedin.com/in/takwa-laffet-883239211}
\quad|\quad
\href{https://github.com/takwa-laffet}
{GitHub: github.com/takwa-laffet}
\quad|\quad
\href{https://laffet-takwa.github.io/portfolio/}
{Portfolio: laffet-takwa.github.io/portfolio}

\end{center}

% =========================================================
% PROFIL
% =========================================================

\section*{Profil professionnel}

Ingénieure en cybersécurité récemment diplômée, avec une expérience pratique
en développement Full-Stack, sécurité applicative, intégration d'API et
automatisation. Maîtrise de React, TypeScript, JavaScript, Python, Node.js,
NestJS, Flask, PostgreSQL, Docker et Linux. Disponible immédiatement pour
un poste de Full-Stack Developer, DevSecOps Engineer ou Software Engineer Junior.

\section*{Formation}

\textbf{ESPRIT -- École Supérieure Privée d'Ingénierie et de Technologies}
\hfill \textbf{2023 -- 2026}

\textit{Diplôme d'Ingénieur en Cybersécurité -- Mention Très Bien / Excellent}

\textbf{ISET Kébili}
\hfill \textbf{2020 -- 2023}

\textit{Licence en Technologie de l'Informatique -- Développement des Systèmes d'Information}


% =========================================================
% EXPERIENCES
% =========================================================

\section*{Expériences professionnelles}

\textbf{CYBEXPSU -- Cybersecurity Center of Excellence, Prince Sultan University,
Arabie Saoudite}
\hfill \textbf{01/2026 -- 06/2026}

\textit{Stagiaire Ingénieure Cybersécurité / Développeuse Full-Stack -- Remote}

\begin{itemize}
    \item Développement d'une plateforme de cybersécurité avec React, Vite,
    Flask et API REST ; création de dashboards d'analyse des alertes SOC.
    \item Intégration de Wazuh, TheHive, Cortex, MISP et Threat Intelligence.
    \item Environnement Linux, Docker, Git et Agile/Scrum.
\end{itemize}

\textbf{Digital Research \& Technologies -- Ariana, Tunisie}
\hfill \textbf{07/2025 -- 09/2025}

\textit{Stagiaire Développeuse Frontend / Full-Stack}

\begin{itemize}
    \item Développement d'une application d'identification et d'analyse de
    contenus avec React, TypeScript, Vite et Tailwind CSS.
    \item Intégration d'API et développement de composants frontend réutilisables.
\end{itemize}

\textbf{Digital Research \& Technologies -- Ariana, Tunisie}
\hfill \textbf{07/2024 -- 09/2024}

\textit{Stagiaire Développeuse Full-Stack}

\begin{itemize}
    \item Développement d'une application de gestion des dons, événements et
    utilisateurs avec NestJS, Prisma, PostgreSQL, Nuxt.js, Vue.js et Vuetify.
    \item Conception d'API REST et fonctionnalités CRUD.
\end{itemize}

\textbf{IT KONCEPT -- Tunisie}
\hfill \textbf{01/2023 -- 06/2023}

\textit{Stagiaire Développeuse Logiciel -- Projet de Fin d'Études}

\begin{itemize}
    \item Développement d'une solution d'automatisation des factures avec
    Odoo 14, Python, PostgreSQL et OCR Mindee.
    \item Développement de modules Odoo avec XML, UML et Scrum.
\end{itemize}

\textbf{AURES -- Tunisie}
\hfill \textbf{01/2022}

\textit{Stagiaire Développeuse Web}

\begin{itemize}
    \item Développement d'une application de suivi du temps avec React,
    Bootstrap et Supabase ; fonctionnalités CRUD.
\end{itemize}

\textbf{Tunisie Télécom -- Tunisie}
\hfill \textbf{2020 -- 2021}

\textit{Stagiaire Réseaux et Infrastructure}

\begin{itemize}
    \item Maintenance et dépannage réseau : TCP/IP, VLAN et configuration réseau.
\end{itemize}

% =========================================================
% COMPETENCES
% =========================================================

\section*{Compétences techniques}

\textbf{Développement :}
Python, JavaScript, TypeScript, SQL, HTML5, CSS3, React.js, Vue.js, Nuxt.js,
Vite, Tailwind CSS, Bootstrap, Node.js, Express.js, NestJS, Flask, FastAPI,
REST API

\textbf{Bases de données :}
PostgreSQL, MySQL, Supabase, Prisma

\textbf{DevOps / DevSecOps :}
Git, GitHub, Docker, Linux, CI/CD, conteneurisation, automatisation,
Secure Software Development, API Security

\textbf{Cybersécurité :}
Wazuh, Suricata, TheHive, Cortex, MISP, n8n, SIEM, Threat Intelligence,
Network Security, Application Security, MITRE ATT\&CK, AWS Cloud Security

% =========================================================
% PROJETS
% =========================================================

\section*{Projets}

\textbf{Plateforme Full-Stack de Gestion des Dons et Événements}
\hfill \textbf{2024}

Nuxt.js, Vue.js, Vuetify, NestJS, Prisma, PostgreSQL. Gestion des utilisateurs,
dons et événements, API REST, CRUD et intégration frontend--backend.

\textbf{Application Web d'Identification et d'Analyse de Contenus}
\hfill \textbf{2025}

React, TypeScript, Vite, Tailwind CSS et API REST. Identification et analyse
de contenus numériques, gestion des résultats et composants réutilisables.

\textbf{Plateforme E-Commerce Full-Stack}
\hfill \textbf{2022--2023}

Développement d'une plateforme e-commerce permettant la gestion des produits,
utilisateurs et commandes. Conception d'interfaces web responsives,
intégration frontend--backend, gestion des données et fonctionnalités CRUD.

% =========================================================
% CERTIFICATIONS
% =========================================================

\section*{Certifications}

Cisco Junior Cybersecurity Analyst \;|\;
Cisco Ethical Hacker \;|\;
Cisco Network Defense \;|\;
AWS Academy Cloud Security Foundations \;|\;
Linux Foundation LFD121

% =========================================================
% LANGUES
% =========================================================

\section*{Langues}

\textbf{Arabe :} langue maternelle \quad
\textbf{Français :} B2 \quad
\textbf{Anglais :} B2

\end{document}